\documentclass[11pt,a4paper,oneside]{report}             % Single-side

%\PassOptionsToPackage{chapternumber=Huordinal}{magyar.ldf}
\usepackage{t1enc}
\usepackage[utf8]{inputenc}
\usepackage{amsmath}
\usepackage{amssymb}
\usepackage{enumerate}
\usepackage[amsmath,thmmarks]{ntheorem}
\usepackage{graphics}
\usepackage{epsfig}
\usepackage{listings}
\usepackage{color}
%\usepackage{fancyhdr}
\usepackage{lastpage}
\usepackage{anysize}
\usepackage[magyar]{babel}
\usepackage{sectsty}
\usepackage{setspace}  % Ettol a tablazatok, abrak, labjegyzetek maradnak 1-es sorkozzel!
\usepackage[hang]{caption}
\usepackage{hyperref}
\usepackage[final]{pdfpages}
%\usepackage{minted}
\usepackage{lscape}
\usepackage{float}
\usepackage{subcaption}
\usepackage{enumitem}
\usepackage{wrapfig}


%--------------------------------------------------------------------------------------
% Main variables
%--------------------------------------------------------------------------------------
\newcommand{\vikszerzo}{Danyi Dávid}
\newcommand{\vikkonzulens}{Kovács Viktor}
\newcommand{\vikcim}{3D rekonstrukció Kinect alkalmazásával}
\newcommand{\viktanszek}{Automatizálási és Alkalmazott Informatika Tanszék}
\newcommand{\vikdoktipus}{Beszámoló}
\newcommand{\vikdepartmentr}{Danyi Dávid}

%--------------------------------------------------------------------------------------
% Page layout setup
%--------------------------------------------------------------------------------------
% we need to redefine the pagestyle plain
% another possibility is to use the body of this command without \fancypagestyle
% and use \pagestyle{fancy} but in that case the special pages
% (like the ToC, the References, and the Chapter pages)remain in plane style

\pagestyle{plain}
\setlength{\parindent}{0pt} % áttekinthetőbb, angol nyelvű dokumentumokban jellemző
\setlength{\parskip}{8pt plus 3pt minus 3pt} % áttekinthetőbb, angol nyelvű dokumentumokban jellemző
%\setlength{\parindent}{12pt} % magyar nyelvű dokumentumokban jellemző
%\setlength{\parskip}{0pt}    % magyar nyelvű dokumentumokban jellemző

\marginsize{35mm}{25mm}{15mm}{15mm} % anysize package
\setcounter{secnumdepth}{0}
\sectionfont{\large\upshape\bfseries}
\setcounter{secnumdepth}{2}
\singlespacing
\frenchspacing

%--------------------------------------------------------------------------------------
%	Setup hyperref package
%--------------------------------------------------------------------------------------
\hypersetup{
    bookmarks=true,            % show bookmarks bar?
    unicode=false,             % non-Latin characters in Acrobat’s bookmarks
    pdftitle={\vikcim},        % title
    pdfauthor={\vikszerzo},    % author
    pdfsubject={\vikdoktipus}, % subject of the document
    pdfcreator={\vikszerzo},   % creator of the document
    pdfproducer={Producer},    % producer of the document
    pdfkeywords={keywords},    % list of keywords
    pdfnewwindow=true,         % links in new window
    colorlinks=true,           % false: boxed links; true: colored links
    linkcolor=black,           % color of internal links
    citecolor=black,           % color of links to bibliography
    filecolor=black,           % color of file links
    urlcolor=black             % color of external links
}

%--------------------------------------------------------------------------------------
% Set up listings
%--------------------------------------------------------------------------------------
\lstset{
	basicstyle=\scriptsize\ttfamily, % print whole listing small
	keywordstyle=\color{black}\bfseries\underbar, % underlined bold black keywords
	identifierstyle=, 					% nothing happens
	commentstyle=\color{white}, % white comments
	stringstyle=\scriptsize\sffamily, 			% typewriter type for strings
	showstringspaces=false,     % no special string spaces
	aboveskip=3pt,
	belowskip=3pt,
	columns=fixed,
	backgroundcolor=\color{lightgray},
} 		
\def\lstlistingname{listing}

%--------------------------------------------------------------------------------------
% Set up minted
%--------------------------------------------------------------------------------------	
%\definecolor{bg}{rgb}{0.95,0.95,0.95}
%\usemintedstyle{emacs}

%--------------------------------------------------------------------------------------
%	Some new commands and declarations
%--------------------------------------------------------------------------------------
\newcommand{\code}[1]{{\upshape\ttfamily\scriptsize\indent #1}}

% define references
\newcommand{\figref}[1]{\ref{fig:#1}.}
\renewcommand{\eqref}[1]{(\ref{eq:#1})}
\newcommand{\listref}[1]{\ref{listing:#1}.}
\newcommand{\sectref}[1]{\ref{sect:#1}}
\newcommand{\tabref}[1]{\ref{tab:#1}.}

\DeclareMathOperator*{\argmax}{arg\,max}
%\DeclareMathOperator*[1]{\floor}{arg\,max}
\DeclareMathOperator{\sign}{sgn}
\DeclareMathOperator{\rot}{rot}
\definecolor{lightgray}{rgb}{0.95,0.95,0.95}

\author{\vikszerzo}
\title{\viktitle}
\includeonly{
	titlepage,%
	introduction,%
	chapter1,%
	chapter2,%
	chapter3,%
	chapter4,%
	conclusion, %
	appendices,%
}

%--------------------------------------------------------------------------------------
%	Setup captions
%--------------------------------------------------------------------------------------
\captionsetup[figure]{
%labelsep=none,
%font={footnotesize,it},
%justification=justified,
width=.75\textwidth,
aboveskip=10pt}

\renewcommand{\captionlabelfont}{\small\bf}
\renewcommand{\captionfont}{\footnotesize\it}

%--------------------------------------------------------------------------------------
% Table of contents and the main text
%--------------------------------------------------------------------------------------
\begin{document}
\pagenumbering{arabic}
\onehalfspacing
%--------------------------------------------------------------------------------------
%	The title page
%--------------------------------------------------------------------------------------
\begin{titlepage}
\begin{center}
\includegraphics[width=60mm,keepaspectratio]{figures/BMElogo.png}\\
\vspace{0.3cm}
%\textbf{Budapest University of Technology and Economics}\\
\textbf{Budapesti Műszaki és Gazdaságtudományi Egyetem}\\
%\textmd{Faculty of Electrical Engineering and Informatics}\\
\textmd{Villamosmérnöki és Informatikai Kar}\\
\textmd{\viktanszek}\\[5cm]

\vspace{0.4cm}
{\huge \bfseries \vikcim}\\[0.8cm]
\vspace{0.5cm}
\textsc{\Large \vikdoktipus}\\[4cm]

\begin{tabular}{cc}
 \makebox[7cm]{\emph{Hallgató}} & \makebox[7cm]{\emph{Konzulens}} \\
 \makebox[7cm]{\vikszerzo} & \makebox[7cm]{\vikkonzulens}
\end{tabular}

\vfill
{\large \today}
\end{center}
\end{titlepage}
\tableofcontents\vfill
%-----------------------------------------------------------------------------------------------
\chapter*{Bevezetés}\addcontentsline{toc}{chapter}{Introduction}
%-----------------------------------------------------------------------------------------------

Ez a dokumentum a 2016 őszi félévében, Önálló laboratórium 2 tárgy keretei között végzett munkám összefoglalója.

Az itt közölt eredmények építenek az előző féléves, azonos témában végzett kutatásomra.
Akkor a feladat a strukturált fényt használó 3D rekonstrukciós eljárások vizsgálata volt.
Az elvek gyakorlati kipróbálására a Microsoft Kinect adott kiváló platformot.
Az előző féléves munka legjavát a technológiával és módszerekkel való ismerkedés adta.
A Kinect által szolgáltatott infravörös kép elemzésével próbáltam reprodukálni az eszköz belső működését.

Az előző félév munkája proof-of-concept jellegű volt.
A mostani ezen túlmutat.
A cél most kettős: egy hosszútávon használható, rugalmas, moduláris keretrendszer fejlesztése a diszparitás meghatározásához, valamint rekonstrukció minőségének javítása a kép lokális struktúrájának figyelembe vételével.

Az első fejezetben röviden összefoglalom a használt algoritmusokat és paraméterezésüket.
Ez részben az előző féléves munka összefoglalása is.

A második, egyben leghosszabb fejezet tartalmazza a fejlesztett keretrendszer leírását.
Ismertetésre kerül a program felhasználó felülete, valami a programozási struktúra és a fejlesztői interfész is.

A harmadik fejezet a lokális struktúra figyelembevételével foglalkozik.
Az itt tárgyalt algoritmusok kísérleti jellegűek, a későbbiekben behatóbb vizsgálatot és optimalizációt igényelnek.

A negyedik részegység az eredmények rövid összegzését és néhány példát tartalmaz.

%-----------------------------------------------------------------------------------------------
\chapter{Development platform}\label{sect:DevTools}
%-----------------------------------------------------------------------------------------------

%-----------------------------------------------------------------------------------------------
\chapter{Teszt platform}\label{sect:DataAnal}
%-----------------------------------------------------------------------------------------------

%-----------------------------------------------------------------------------------------------
\chapter{Lokális struktúra vizsgálata}\label{sect:LocalStruct}
%-----------------------------------------------------------------------------------------------

Az illeszkedési minőség bizonyos esetekben javítható, ha figyelembe vesszük a kép lokális struktúráját.
Ez a megoldás rontja az általánosságot, mert extra megkötéseket követel meg a diszparitásra vonatkozóan.

Az alábbiakban az ilyen feltételeket és a rájuk illesztett modelleket ismertetem.
Nem minden modell került megvalósításra, néhol a vizsgálati eredményekből kiderült, nem reális vagy fölösleges részletes.

%-----------------------------------------------------------------------------------------------
\section{Modellek}\label{sect:Models}
%-----------------------------------------------------------------------------------------------

Ezen fejezet alapfeltevése, hogy a diszparitás kis régiókon belül nem változik tetszőleges.
Kis régió fogalom nagyon szubjektív, jelen állás szerint kísérletileg meghatározott fogalom.
A továbbiak a kép 16 pixel széles tartományát értem ezalatt.

Már korábban is feltettük\footnote{Az ablakméret megválasztásánál}, hogy nem számítunk tetszőlegesen nagy elmozdulásokra.
A modellezés során ennél erősebb feltételek keresése volt a cél.

Két esetet vizsgáltam a féléves munka során: lineáris modell és 2 konstans szakaszból álló modell.

%-----------------------------------------------------------------------------------------------
\subsection{Lineáris modell}\label{sect:linearModel}
%-----------------------------------------------------------------------------------------------

A lineáris modell esetében az volt a feltételezés, hogy diszparitás lineárisan változik rövid szakaszon.
Ekkor a modell \eqref{linModel} alakú.
\begin{equation}
y = a*x + b
\label{eq:linModel}
\end{equation}

Az illeszkedés hibája nyilvánvalóan igen nagy lesz, ha a diszparitásban ugrás\footnote{Ez a helyzet áll fenn például tárgyak pereménél} van a régión belül.
Kiegészíthető a modell úgy, hogy 2 egyenes szakaszt illesztünk az adatokra.
Így jóval több paraméter\footnote{Az egyenesek paraméterein túl a töréspont helyét is meg kell határozni.} meghatározása a feladat ugyan annyi bemenő adatpont használatával.
Ha ezt a problémát a régióméret növelésével próbáljuk csökkenteni, akkor pedig a modell alapfeltevése (lineáris változás) sérülhet.

Mélyebb vizsgálatnak nem vetettem alá ezt a modellt.
A jelenlegi pontossági igényeken belül még a lineáris változás is elhanyagolható kis szakaszokon.
Több mintát megvizsgálva arra jutottam, hogy a konstans modell is elegendő a feladat megoldására.

%-----------------------------------------------------------------------------------------------
\subsection{Konstans modell}\label{sect:constModel}
%-----------------------------------------------------------------------------------------------

Hasonlóan a lineáris modellhez, ebben az esetben is be meg kell engedni egy töréspontot modellben.
Így egy három paraméteres modell adódik:
\begin{itemize}[noitemsep]
\item Diszparitás a régió elején
\item Diszparitás a régió végén
\item Törés koordinátája
\end{itemize}

Még így is relatív kevés adat jut egy paraméterre (16 pixeles régióméret esetén).
A gyakorlati implementációban további optimalizálásként 2 esetet különböztetek meg.
Ha a két diszparitás paraméter eltérése igen kicsi (1-2 pixel), akkor 1 paraméteres modellt használok.

%-----------------------------------------------------------------------------------------------
\section{Illesztési módszerek}\label{sect:modelMatch}
%-----------------------------------------------------------------------------------------------

Az előző féléves munkám során a rekonstrukció eredményeként mindig a legjobb illeszkedési mutatóval rendelkező pont lett elfogadva.
A mostani módszerek ettől gyökeresen különböznek.
Az első maximumok gyakran hibás egyezés következményei, zajosak lehetnek.
Az vizualizált illeszkedési gráfokat vizsgálva szembetűnik, hogy sok esetben az első maximum ugyan hibás, de a helyes eredmény is detektálható lenne az adatokból.
Itt jön a képbe ezen fejezet tárgya: a lokális környezet.
Ha nem csak az első maximumot tároljuk, hanem az első néhányat, a szomszédos pontok diszparitása alapján kiválaszthatjuk a legvalószínűbb egyezést közülük.

Ezt a választást segítik az feljebb vázolt modellek.
Alább pedig két lehetőséget mutatok be paraméterek meghatározására.

%-----------------------------------------------------------------------------------------------
\subsection{Kumulált korreláció}\label{sect:CumCorr}
%-----------------------------------------------------------------------------------------------



%-----------------------------------------------------------------------------------------------
\subsection{RANSAC}\label{sect:ransac}
%-----------------------------------------------------------------------------------------------
%-----------------------------------------------------------------------------------------------
\chapter{Eredmények}\label{sect:Results}
%-----------------------------------------------------------------------------------------------

%-----------------------------------------------------------------------------------------------
\chapter{Összegzés}\label{sect:conclusion}
%-----------------------------------------------------------------------------------------------


%\listoffigures\addcontentsline{toc}{chapter}{Ábrák jegyzéke}
%\listoftables\addcontentsline{toc}{chapter}{Táblázatok jegyzéke}

\bibliography{mybib}
\addcontentsline{toc}{chapter}{Bibliography}
\bibliographystyle{plain}

%----------------------------------------------------------------------------
\appendix
%----------------------------------------------------------------------------


\label{page:last}
\end{document}